\begin{landscape}
\begingroup
\fontsize{10}{13}\selectfont
\fancyhead[L]{\small POSTHUMOUS PAPERS. PERTURBATIONS OF PALLAS}
\renewcommand{\headrulewidth}{0pt}
\par\medskip
\pdfbookmark[2]{Posthumous papers VI [10]}{pallas-VI-fragment10}
\begin{center}[10.]\end{center}
\par\noindent\hspace*{\parindent}[For the exact integration, set
\[
u=18\text{\mdseries\AstroSymbols ♃}-7\text{\mdseries\AstroSymbols ⚴}+A,\qquad
A=\text{constant},\qquad
\int n\mathrm{d}t=\alpha\iint\sin u\,\mathrm{d}\text{\mdseries\AstroSymbols ⚴}^{2}.
\]
\SourceLeaf{564}
\par\noindent\hspace*{\parindent}Now since
\[
\text{\mdseries\AstroSymbols ⚴}=\int n\mathrm{d}t+\varepsilon,\quad
\text{hence}\quad\frac{\mathrm{d}\mathrm{d}\text{\mdseries\AstroSymbols ⚴}}{\mathrm{d}t^2}=\frac{\mathrm{d}n}{\mathrm{d}t},
\]
we obtain the equation for \(u\)
\[
\frac{\mathrm{d}\mathrm{d}u}{nn\mathrm{d}t^2}=-7\alpha\sin u,
\]
hence
\[
\left(\frac{\mathrm{d}u}{n\mathrm{d}t}\right)^2=C+14\alpha\cos u.
\]
\par\noindent\hspace*{\parindent}It follows that if \(C\) is positive and greater in magnitude than \(14\alpha\), then \(\dfrac{\mathrm{d}u}{\mathrm{d}t}\) does not vanish and \(u\) increases or decreases continuously. If \(C\) is positive and smaller than the magnitude of \(14\alpha\), or if \(C\) is negative, then \(u\) is a periodic function and oscillates about \(0\) or \(180^0\) according as \(\alpha\) is positive or negative. In the latter case, the mean motions are thus in a rational ratio and the ratio \(\dfrac{n'}{n}=\tfrac{7}{18}\) is repeatedly restored.
\par\noindent\hspace*{\parindent}If we denote the value of \(u\) at an arbitrary initial time by \(U\) and the corresponding value of \(\dfrac{\mathrm{d}u}{\mathrm{d}\text{\mdseries\AstroSymbols ⚴}}=18\dfrac{n'}{n}-7\) by \(\mu\), we have
\[
C=\mu\mu-14\alpha\cos U
\]
and with the numerical values
\begin{center}
\begin{tabular}{@{}lll@{}}
\(\alpha=0\overset{\prime\prime}{,}02955\)&\(n'=299\overset{\prime\prime}{,}12817\)&\(18\text{\mdseries\AstroSymbols ♃}-7\text{\mdseries\AstroSymbols ⚴}=125^0 36'\,49''\)\\
\(A=80^0 56'\,57''\)&\(n=769{,}16512\)&\(U=206\quad33\quad46\)
\end{tabular}
\end{center}
we obtain
\[
\begin{gathered}
\log14\alpha=4{,}30221-10\\
\log\mu=6{,}29314-10
\end{gathered}
\qquad\log C=4{,}26300-10.
\]
Thus this is indeed a case in which the mean motions are in a rational ratio.
\par\noindent\hspace*{\parindent}Furthermore,
\[
\left(\frac{\mathrm{d}u}{n\mathrm{d}t}\right)^2=(C+14\alpha)\left\{1-\frac{28\alpha}{C+14\alpha}\sin\tfrac12u^2\right\},
\]
and by the substitution
\[
\BareRoot{\frac{28\alpha}{C+14\alpha}}\mathpunct{\cdot}\sin\tfrac12u=\sin\psi
\]
\[
\frac{\mathrm{d}u}{n\mathrm{d}t}=\BareRoot{(C+14\alpha)}\mathpunct{.}\cos\psi
\]
we obtain
\[
\frac{\mathrm{d}\psi}{n\mathrm{d}t}=\BareRoot{7\alpha}\mathpunct{.}\BareRoot{\left(1-\frac{C+14\alpha}{28\alpha}\sin\psi^2\right)}
\]
hence
\[
n\mathrm{d}t=\frac{\mathrm{d}\psi}{\BareRoot{\left(7\alpha-\dfrac{C+14\alpha}{4}\sin\psi^2\right)}}\mathpunct{\cdot}
\]
\par\noindent\hspace*{\parindent}The maximum (respectively, minimum) of \(\dfrac{\mathrm{d}u}{n\mathrm{d}t}\) is
\[
=\BareRoot{(C+14\alpha)}=\BareRoot{\left(28\alpha\sin\frac{U^2}{2}+\mu\mu\right)}.
\]
\SourceLeaf{565}
\par\noindent\hspace*{\parindent}Accordingly, the mean motion of Pallas oscillates between the values:
\[
\frac{18}{7}n'\pm0\overset{\prime\prime}{,}2153.
\]
\par\noindent\hspace*{\parindent}The ratio of the period of this oscillation to one revolution of Pallas is given by the integral, taken from \(0\) to \(\dfrac{\pi}{2}\):
\[
\begin{aligned}
\frac{2}{\pi}\int\frac{\mathrm{d}\psi}{\BareRoot{\left(7\alpha-\dfrac{C+14\alpha}{4}\sin\psi^2\right)}}
&=\frac{1}{M\left(\BareRoot{7\alpha},\BareRoot{\dfrac{14\alpha-C}{4}}\right)}\\[.8ex]
&=\frac{1}{M(\BareRoot{7\alpha},\BareRoot{(7\alpha\cos\tfrac12U^2-\tfrac14\mu\mu)})}\mathpunct{\raisebox{.35ex}{,}}
\end{aligned}
\]
that is, 1894 revolutions of Pallas \(=\) 737 revolutions of Jupiter.
\par\noindent\hspace*{\parindent}The maximum and minimum of \(u\) itself are:
\[
\pm156^0 1'.\text{]}
\]
\clearpage
\endgroup
\end{landscape}
